\lyseis{Β2.4}{Συμμετρία ως προς σημείο}
\sol{1} Ελεύθερη κατασκευή. Στην προέκταση της $\mathrm{ΑΟ}$ παίρνω $\mathrm{ΟΑ}'=\mathrm{ΟΑ}$.
\sol{2} $\mathrm{ΟΑ}'=3{,}2$ cm και $\mathrm{ΑΑ}'=6{,}4$ cm.
\sol{3} Το συμμετρικό τμήμα $\mathrm{Α}'\mathrm{Β}'$ είναι ίσο και παράλληλο με το $\mathrm{ΑΒ}$.
\sol{4} Ελεύθερη κατασκευή. Το συμμετρικό τρίγωνο είναι ίσο και «ανάποδα στραμμένο».
\sol{5} $-3{,}5$, $2$ και $0$.
\sol{6} Ως προς το $1$: το συμμετρικό του $3{,}5$ είναι $-1{,}5$. Ως προς το $3$: το συμμετρικό του $-2$ είναι $8$ (αφού $\frac{-2+8}{2}=3$).
\sol{7} α) Σωστό. β) Σωστό. γ) Σωστό. δ) Λάθος: είναι παράλληλο.
\sol{8} Στην απέναντι θέση, $12$ μ. από το κέντρο στην προέκταση. Οι δύο θέσεις απέχουν $24$ μ.
\sol{9} Στη συμμετρική θέση ως προς το $\mathrm{Ο}$, δηλαδή στην προέκταση της ευθείας φυτό--$\mathrm{Ο}$, πάλι $2{,}5$ μ. από το $\mathrm{Ο}$.
\sol{10} Δείχνει προς τα κάτω. Οι δύο θέσεις της άκρης είναι συμμετρικές ως προς το κέντρο του ρολογιού.
\sol{11} Στην αντίθετη θέση, γιατί οι δύο πτερύγια είναι συμμετρικά ως προς το κέντρο.
\sol{12} Η λέξη φαίνεται ανάποδα (αναποδογυρισμένη), αλλά όχι αντικατοπτρισμένη.
\sol{13} Ως προς σημείο: το συμμετρικό τρίγωνο προκύπτει με στροφή $180^\circ$ και δεν αναποδογυρίζει σαν καθρέφτης. Ως προς άξονα το συμμετρικό είναι «καθρεφτισμένο».
\sol{14} Το $\mathrm{Α}$.
\sol{15} Το $\mathrm{ΑΒΑ}'\mathrm{Β}'$ έχει διαγώνιους $\mathrm{ΑΑ}'$ και $\mathrm{ΒΒ}'$ που διχοτομούνται στο $\mathrm{Ο}$: είναι παραλληλόγραμμο, άρα $\mathrm{ΑΒ}\parallel\mathrm{Α}'\mathrm{Β}'$.
\sol{16} Κάθε ευθεία που περνά από το $\mathrm{Ο}$.
\sol{17} Ελεύθερη κατασκευή.
\sol{18} Ελεύθερη απάντηση.
